% GENERATED FILE. Edit canonical lessons, then run npm run generate:books.
% canonical-source-hash: fnv1a64:7a1bdf94c50fd954
% canonical-lessons: MR-C533-kasrat, MR-C533-yogasan, MR-C533-sitti, MR-C533-vakne, MR-C533-jhepavne

\chapter{Keeping Fit}
\label{ch:mr-keeping-fit}
\hlchaptermodality{533}

\begin{chapteropening}
\textbf{By the end of this chapter:} \emph{I can talk about workouts, yoga and moving the body.}

Complete the last lesson of chapter 533: I can talk about workouts, yoga and moving the body.
\end{chapteropening}

\section[\texorpdfstring{\mr{कसरत} (kasrat)}{कसरत (kasrat)}]{\mr{कसरत} (kasrat) — a workout; exercise}

\label{lesson:MR-C533-kasrat}



\begin{quote}
Before the new one: say the Marathi for a crossword, then the Marathi for hide and seek.
\end{quote}

\subsection*{You'll want to know: \mr{कसरत}}
\textbf{\mr{कसरत}} — \emph{kasrat} — \textquotedblleft{}a workout; exercise\textquotedblright{}.

\begin{quote}
\mr{धावपटू} \mr{रोज} \mr{दोन} \mr{तास} \mr{कसरत} \mr{करतात}.
\end{quote}

\begin{quote}
\emph{dhāvpaṭū roj don tās kasrat kartāt.}
\end{quote}

\begin{quote}
Runners work out for two hours every day.
\end{quote}

Feminine.

\mr{कसरत} \mr{करणे}, to work out.

\begin{cousinweb}
\textbf{\mr{कसरत}} (\emph{kasrat}) is usually traced to Persian \emph{kasrat}, \textquotedblleft{}exercise\textquotedblright{}.
\end{cousinweb}

\begin{grammarlens}[title={What you've built}]
One more word for describing what you lived and felt.
\end{grammarlens}

\subsection*{Your turn}
\begin{itemize}
\raggedright
  \item \emph{Say it:} \emph{kasrat}
  \item \emph{Say it:} \emph{kasrat}, once more
  \item \emph{Say it:} \emph{dhāvpaṭū roj don tās kasrat kartāt.} — Runners work out for two hours every day.
  \item \emph{Recall:} say the Marathi for a crossword, then the Marathi for hide and seek, then say \emph{kasrat} again
\end{itemize}

\begin{tcolorbox}[breakable,colback=teal!4,colframe=teal!35!black,title={Before you move on}]
What does \mr{कसरत} mean? (\textquotedblleft{}A workout; exercise\textquotedblright{}.) Say it once more.
\end{tcolorbox}
\section[\texorpdfstring{\mr{योगासन} (yogāsan)}{योगासन (yogāsan)}]{\mr{योगासन} (yogāsan) — a yoga posture}

\label{lesson:MR-C533-yogasan}



\begin{quote}
Before the new one: say the Marathi for hide and seek, then the Marathi for a workout.
\end{quote}

\subsection*{You'll want to know: \mr{योगासन}}
\textbf{\mr{योगासन}} — \emph{yogāsan} — \textquotedblleft{}a yoga posture\textquotedblright{}.

\begin{quote}
\mr{आजी} \mr{रोज} \mr{सकाळी} \mr{योगासनं} \mr{करते}.
\end{quote}

\begin{quote}
\emph{ājī roj sakāḷī yogāsanaṁ karte.}
\end{quote}

\begin{quote}
Grandmother does yoga every morning.
\end{quote}

Neuter. Plural: \mr{योगासनं}, \emph{yogāsanaṁ}.

\begin{cousinweb}
\textbf{\mr{योगासन}} (\emph{yogāsan}) is a tatsama compound of Sanskrit \emph{yoga} and \emph{āsana}, \textquotedblleft{}a seat, a posture\textquotedblright{} (Monier-Williams).
\end{cousinweb}

\begin{grammarlens}[title={What you've built}]
One more word for describing what you lived and felt.
\end{grammarlens}

\subsection*{Your turn}
\begin{itemize}
\raggedright
  \item \emph{Say it:} \emph{yogāsan}
  \item \emph{Say it:} \emph{yogāsan}, once more
  \item \emph{Say it:} \emph{ājī roj sakāḷī yogāsanaṁ karte.} — Grandmother does yoga every morning.
  \item \emph{Recall:} say the Marathi for hide and seek, then the Marathi for a workout, then say \emph{yogāsan} again
\end{itemize}

\begin{tcolorbox}[breakable,colback=teal!4,colframe=teal!35!black,title={Before you move on}]
What does \mr{योगासन} mean? (\textquotedblleft{}A yoga posture\textquotedblright{}.) Say it once more.
\end{tcolorbox}
\section[\texorpdfstring{\mr{शिट्टी} (śiṭṭī)}{शिट्टी (śiṭṭī)}]{\mr{शिट्टी} (śiṭṭī) — a whistle}

\label{lesson:MR-C533-sitti}



\begin{quote}
Before the new one: say the Marathi for a workout, then the Marathi for a yoga posture.
\end{quote}

\subsection*{You'll want to know: \mr{शिट्टी}}
\textbf{\mr{शिट्टी}} — \emph{śiṭṭī} — \textquotedblleft{}a whistle\textquotedblright{}.

\begin{quote}
\mr{प्रशिक्षकाने} \mr{शिट्टी} \mr{वाजवली} \mr{आणि} \mr{खेळ} \mr{सुरू} \mr{झाला}.
\end{quote}

\begin{quote}
\emph{praśikṣakāne śiṭṭī vājavlī āṇi kheḷ surū jhālā.}
\end{quote}

\begin{quote}
The coach blew the whistle and the game began.
\end{quote}

Feminine.

\mr{शिट्टी} \mr{वाजवणे}, to blow a whistle.

\begin{cousinweb}
\textbf{\mr{शिट्टी}} (\emph{śiṭṭī}) is inherited: CDIAL lists Marathi \emph{śiṭī}, \textquotedblleft{}whistle\textquotedblright{}, under the sound word \emph{*sīṭṭa} (CDIAL 13427).
\end{cousinweb}

\begin{grammarlens}[title={What you've built}]
One more word for describing what you lived and felt.
\end{grammarlens}

\subsection*{Your turn}
\begin{itemize}
\raggedright
  \item \emph{Say it:} \emph{śiṭṭī}
  \item \emph{Say it:} \emph{śiṭṭī}, once more
  \item \emph{Say it:} \emph{praśikṣakāne śiṭṭī vājavlī āṇi kheḷ surū jhālā.} — The coach blew the whistle and the game began.
  \item \emph{Recall:} say the Marathi for a workout, then the Marathi for a yoga posture, then say \emph{śiṭṭī} again
\end{itemize}

\begin{tcolorbox}[breakable,colback=teal!4,colframe=teal!35!black,title={Before you move on}]
What does \mr{शिट्टी} mean? (\textquotedblleft{}A whistle\textquotedblright{}.) Say it once more.
\end{tcolorbox}
\section[\texorpdfstring{\mr{वाकणे} (vākṇe)}{वाकणे (vākṇe)}]{\mr{वाकणे} (vākṇe) — to bend}

\label{lesson:MR-C533-vakne}



\begin{quote}
Before the new one: say the Marathi for a yoga posture, then the Marathi for a whistle.
\end{quote}

\subsection*{You'll want to know: \mr{वाकणे}}
\textbf{\mr{वाकणे}} — \emph{vākṇe} — \textquotedblleft{}to bend\textquotedblright{}.

\begin{quote}
\mr{चेंडू} \mr{उचलण्यासाठी} \mr{तो} \mr{खाली} \mr{वाकला}.
\end{quote}

\begin{quote}
\emph{ceṇḍū ucalṇyāsāṭhī to khālī vāklā.}
\end{quote}

\begin{quote}
He bent down to pick up the ball.
\end{quote}

Past (he …): \mr{तो} \mr{वाकला}, \emph{to vāklā}.

\begin{cousinweb}
\textbf{\mr{वाकणे}} (\emph{vākṇe}) is built on \mr{वाक}, \textquotedblleft{}a curve\textquotedblright{}, which CDIAL sets under Sanskrit \emph{vaṅka}, \textquotedblleft{}bent, crooked\textquotedblright{} (CDIAL 11191).
\end{cousinweb}

\begin{grammarlens}[title={What you've built}]
One more word for describing what you lived and felt.
\end{grammarlens}

\subsection*{Your turn}
\begin{itemize}
\raggedright
  \item \emph{Say it:} \emph{vākṇe}
  \item \emph{Say it:} \emph{vākṇe}, once more
  \item \emph{Say it:} \emph{ceṇḍū ucalṇyāsāṭhī to khālī vāklā.} — He bent down to pick up the ball.
  \item \emph{Recall:} say the Marathi for a yoga posture, then the Marathi for a whistle, then say \emph{vākṇe} again
\end{itemize}

\begin{tcolorbox}[breakable,colback=teal!4,colframe=teal!35!black,title={Before you move on}]
What does \mr{वाकणे} mean? (\textquotedblleft{}To bend\textquotedblright{}.) Say it once more.
\end{tcolorbox}
\section[\texorpdfstring{\mr{झेपावणे} (jhepāvṇe)}{झेपावणे (jhepāvṇe)}]{\mr{झेपावणे} (jhepāvṇe) — to leap; to spring}

\label{lesson:MR-C533-jhepavne}



\begin{quote}
Before the new one: say the Marathi for a whistle, then the Marathi for to bend.
\end{quote}

\subsection*{You'll want to know: \mr{झेपावणे}}
\textbf{\mr{झेपावणे}} — \emph{jhepāvṇe} — \textquotedblleft{}to leap; to spring\textquotedblright{}.

\begin{quote}
\mr{चेंडू} \mr{पकडण्यासाठी} \mr{तो} \mr{पुढे} \mr{झेपावला}.
\end{quote}

\begin{quote}
\emph{ceṇḍū pakaḍṇyāsāṭhī to puḍhe jhepāvlā.}
\end{quote}

\begin{quote}
He sprang forward to catch the ball.
\end{quote}

Past (he …): \mr{तो} \mr{झेपावला}, \emph{to jhepāvlā}.

Built on \mr{झेप}, a leap.

\begin{cousinweb}
\textbf{\mr{झेपावणे}} (\emph{jhepāvṇe}) is built on \mr{झेप}, \textquotedblleft{}a leap\textquotedblright{}; its origin is not settled.
\end{cousinweb}

\begin{grammarlens}[title={What you've built}]
One more word for describing what you lived and felt.
\end{grammarlens}

\subsection*{Your turn}
\begin{itemize}
\raggedright
  \item \emph{Say it:} \emph{jhepāvṇe}
  \item \emph{Say it:} \emph{jhepāvṇe}, once more
  \item \emph{Say it:} \emph{ceṇḍū pakaḍṇyāsāṭhī to puḍhe jhepāvlā.} — He sprang forward to catch the ball.
  \item \emph{Recall:} say the Marathi for a whistle, then the Marathi for to bend, then say \emph{jhepāvṇe} again
\end{itemize}

\begin{tcolorbox}[breakable,colback=teal!4,colframe=teal!35!black,title={Before you move on}]
What does \mr{झेपावणे} mean? (\textquotedblleft{}To leap; to spring\textquotedblright{}.) Say it once more.
\end{tcolorbox}
